\MFQTransition{组合实施、应用分析与开放项目 题组}
\begin{enumerate}[leftmargin=*]
  \item \textbf{口述概念：}信号、排序、目标权重、成交权重和净收益分别是什么对象？
  \item \textbf{口述概念：}等权、市值权重与重叠持有期改变了什么研究问题？
  \item \textbf{笔试推导：}手算两期等权组合的毛收益、换手、成本和净收益。
  \item \textbf{笔试推导：}写出 $H$ 期重叠持有的有效权重及其推断后果。
  \item \textbf{数值编程：}实现五分组、等权/市值权重与 1/3/6 期持有敏感性表。
  \item \textbf{数值编程：}从各期实际持仓和下一期资产收益计算组合毛收益，再根据成交增量扣除成本。
  \item \textbf{研究判断：}审计“只保留 20 个候选再做 BH”的研究报告。
  \item \textbf{研究判断：}审计未记录停牌、融券与成交完成率的 A 股组合。
  \item \textbf{综合项目：}选择一个因子，完成从回归到分组组合的分析。比较两种持有期或权重规则，并解释成本与样本选择怎样影响结果。
\end{enumerate}

\subsection*{基础衔接：独立计算}
同一组合的总权重与绝对权重和分别是多少？能否把 0.5\% 毛收益称作未加杠杆收益？
